% Chapter 9 — the graph in numbers: the live build (KG 7.0.0), measured.

\gmzoomin{numbers}
\note{Chapter 9: step back and look at the whole graph.}

\begin{frame}{MemeAtlas next to IMKG}
  \vskip2mm
  {\small
  \begin{tabular}{@{}l@{\hspace{8mm}}r@{\hspace{8mm}}r@{}}
    & \bfseries\color{gm@col@numbers}MemeAtlas \NVersion & \bfseries\color{gmMuted}IMKG, KYM part\\\midrule
    Know Your Meme entries        & \textbf{\NFrames}   & \FImkgFrames\\
    nodes (core)                  & \textbf{\NCoreNodes}& \FImkgKymNodes\\
    edges (core)                  & \textbf{\NCoreEdges}& \FImkgKymEdges\\
    \quad counted the RDF way     & \textbf{\NCoreTE}   & \FImkgKymEdges\\
    whole graph: nodes            & \textbf{\NNodes}    & ---\\
    whole graph: edges            & \textbf{\NEdges}    & ---\\
    rebuilt                       & every month         & once\\
  \end{tabular}}
  \vskip6mm
  {\small\color{gmInk2}\NFramesRatio$\times$ the entries --- and every number on this slide is recomputed with each build.}
  \note{
  \begin{itemize}
    \item The fair comparison is with the KYM part of IMKG (its Table 2), the same source: frames,
      entry types, tags, series and links --- the ``core''.
    \item MemeAtlas covers 24,291 entries against 12,585: 1.93 times as many.
    \item Edges: IMKG is RDF, where every literal is an edge; MemeAtlas's property graph keeps
      literals on the nodes. Counted the RDF way (``triple-equivalent''), the core has 921,120
      edges against IMKG's 914,941 --- the same density over twice the frames.
    \item And beyond the core, the layers IMKG did not have: events, templates, curated links.
      The whole graph: 795,711 nodes, 2.56 million edges.
  \end{itemize}}
\end{frame}

\begin{frame}{What each entry carries}
  \begin{columns}[c]
    \column{.56\textwidth}
      {\scriptsize\color{gmInk2}entries with \dots}\par\vskip1mm
      % generated by scripts/figures.py — do not edit
\begin{tikzpicture}
\begin{axis}[gm hbar, scale only axis, width=.8\linewidth, height=3.0cm, xmax=29009, symbolic y coords={{a Wikidata link},{events},{an imgflip template},{all three},{none of them}},
  point meta=explicit symbolic, nodes near coords={\pgfplotspointmeta}]
\addplot[fill=gm@col@numbers] coordinates {(24174,{a Wikidata link}) [99.5\%] (18773,{events}) [77.3\%] (8368,{an imgflip template}) [34.4\%] (8158,{all three}) [33.6\%] (65,{none of them}) [0.3\%]};
\end{axis}
\end{tikzpicture}

    \column{.4\textwidth}
      \gmbig{\NWithNone}{entries with none of the three}\par\vskip3mm
      {\small\color{gmInk2}one in three has everything: links, a story and a picture}
  \end{columns}
  \note{
  \begin{itemize}
    \item 99.5\% of entries have at least one Wikidata link; 77.3\% have events (nearly all that
      have a story); 34.4\% have a template --- templates exist for picture memes, not for people
      or events; 33.6\% have all three.
    \item Only 65 entries have none of them.
  \end{itemize}}
\end{frame}

\begin{frame}{When memes are born}
  {\scriptsize\color{gmInk2}entries by the year their meme started}\par\vskip1mm
  % generated by scripts/figures.py — do not edit
\begin{tikzpicture}
\begin{axis}[gm axis, ybar, bar width=4.2pt, width=\linewidth, height=4.6cm,
  xmin=1994.2, xmax=2026.8, ymin=0, ymax=2130, axis x line*=bottom, axis y line=none,
  xtick={1995,2000,2005,2010,2015,2020,2025}, x tick label style={/pgf/number format/1000 sep={}},
  every axis plot/.append style={draw=none}, clip=false]
\fill[gmHair, opacity=.55] (axis cs:2022.5,0) rectangle (axis cs:2026.5,2092);
\node[font=\tiny, text=gmInk2, anchor=north east, align=right] at (axis cs:2026.4,2055) {KYM writes up\\a meme later};
\addplot[fill=gm@col@numbers] coordinates {(1995,65) (1996,75) (1997,77) (1998,93) (1999,165) (2000,133) (2001,146) (2002,162) (2003,246) (2004,246) (2005,296) (2006,433) (2007,428) (2008,478) (2009,538) (2010,707) (2011,881) (2012,839) (2013,903) (2014,1013) (2015,1013) (2016,1199) (2017,1466) (2018,1549) (2019,1868) (2020,1697) (2021,1347) (2022,1326) (2023,1267) (2024,1067) (2025,577) (2026,251)};
\node[font=\tiny, text=gmInk, anchor=south] at (axis cs:2019,1868) {1,868};
\end{axis}
\end{tikzpicture}

  \vskip2mm
  {\small\color{gmInk2}The peak is 2019. The drop after is mostly time: Know Your Meme documents a meme once it has lasted.}
  \note{
  \begin{itemize}
    \item Every entry has the year its meme started (from KYM's infobox). The curve climbs from the
      late 1990s to a peak in 2019 (1,868 entries).
    \item Careful with the right side: KYM writes a meme up after it has proved itself, so the
      last years are not complete. This is a documentation lag, not fewer memes.
    \item 970 entries start before 1995: KYM also documents older cultural references.
  \end{itemize}}
\end{frame}

\begin{frame}{Where memes are born}
  {\scriptsize\color{gmInk2}entries per year, by the platform their meme started on --- one scale for all six}\par\vskip2mm
  \resizebox{\linewidth}{!}{% generated by scripts/figures.py — do not edit
\begin{tikzpicture}
\begin{scope}[xshift=0.00cm, yshift=0.00cm]
\begin{axis}[gm axis, width=3.9cm, height=1.75cm, scale only axis, at={(0,0)},
  xmin=2005, xmax=2025, ymin=0, ymax=496, axis x line*=bottom, axis y line=none,
  xtick={2005,2015,2025}, x tick label style={/pgf/number format/1000 sep={}, font=\small},
  title={4chan}, title style={font=\large\bfseries, yshift=-1mm}]
\addplot[fill=gm@col@numbers!45!gmSurface, draw=gm@col@numbers, line width=1pt] coordinates {(2005,18) (2006,44) (2007,42) (2008,42) (2009,59) (2010,89) (2011,37) (2012,51) (2013,38) (2014,45) (2015,59) (2016,50) (2017,40) (2018,60) (2019,31) (2020,43) (2021,19) (2022,19) (2023,13) (2024,5) (2025,1)} \closedcycle;
\node[font=\small\bfseries, text=gmInk, anchor=south] at (axis cs:2010,89) {2010};
\end{axis}
\node[font=\small, text=gmInk2, anchor=north west] at (0,-.5) {805 memes};
\end{scope}
\begin{scope}[xshift=4.75cm, yshift=0.00cm]
\begin{axis}[gm axis, width=3.9cm, height=1.75cm, scale only axis, at={(0,0)},
  xmin=2005, xmax=2025, ymin=0, ymax=496, axis x line*=bottom, axis y line=none,
  xtick={2005,2015,2025}, x tick label style={/pgf/number format/1000 sep={}, font=\small},
  title={YouTube}, title style={font=\large\bfseries, yshift=-1mm}]
\addplot[fill=gm@col@numbers!45!gmSurface, draw=gm@col@numbers, line width=1pt] coordinates {(2005,10) (2006,86) (2007,103) (2008,76) (2009,96) (2010,109) (2011,125) (2012,96) (2013,107) (2014,107) (2015,113) (2016,134) (2017,125) (2018,140) (2019,185) (2020,152) (2021,122) (2022,108) (2023,127) (2024,91) (2025,36)} \closedcycle;
\node[font=\small\bfseries, text=gmInk, anchor=south] at (axis cs:2019,185) {2019};
\end{axis}
\node[font=\small, text=gmInk2, anchor=north west] at (0,-.5) {2,248 memes};
\end{scope}
\begin{scope}[xshift=9.50cm, yshift=0.00cm]
\begin{axis}[gm axis, width=3.9cm, height=1.75cm, scale only axis, at={(0,0)},
  xmin=2005, xmax=2025, ymin=0, ymax=496, axis x line*=bottom, axis y line=none,
  xtick={2005,2015,2025}, x tick label style={/pgf/number format/1000 sep={}, font=\small},
  title={Tumblr}, title style={font=\large\bfseries, yshift=-1mm}]
\addplot[fill=gm@col@numbers!45!gmSurface, draw=gm@col@numbers, line width=1pt] coordinates {(2005,1) (2006,4) (2007,0) (2008,5) (2009,14) (2010,29) (2011,85) (2012,67) (2013,87) (2014,102) (2015,77) (2016,69) (2017,67) (2018,50) (2019,24) (2020,20) (2021,27) (2022,22) (2023,16) (2024,6) (2025,4)} \closedcycle;
\node[font=\small\bfseries, text=gmInk, anchor=south] at (axis cs:2014,102) {2014};
\end{axis}
\node[font=\small, text=gmInk2, anchor=north west] at (0,-.5) {776 memes};
\end{scope}
\begin{scope}[xshift=0.00cm, yshift=-3.25cm]
\begin{axis}[gm axis, width=3.9cm, height=1.75cm, scale only axis, at={(0,0)},
  xmin=2005, xmax=2025, ymin=0, ymax=496, axis x line*=bottom, axis y line=none,
  xtick={2005,2015,2025}, x tick label style={/pgf/number format/1000 sep={}, font=\small},
  title={Reddit}, title style={font=\large\bfseries, yshift=-1mm}]
\addplot[fill=gm@col@numbers!45!gmSurface, draw=gm@col@numbers, line width=1pt] coordinates {(2005,1) (2006,3) (2007,2) (2008,9) (2009,13) (2010,32) (2011,94) (2012,96) (2013,69) (2014,50) (2015,47) (2016,56) (2017,104) (2018,174) (2019,212) (2020,121) (2021,81) (2022,69) (2023,49) (2024,37) (2025,13)} \closedcycle;
\node[font=\small\bfseries, text=gmInk, anchor=south] at (axis cs:2019,212) {2019};
\end{axis}
\node[font=\small, text=gmInk2, anchor=north west] at (0,-.5) {1,332 memes};
\end{scope}
\begin{scope}[xshift=4.75cm, yshift=-3.25cm]
\begin{axis}[gm axis, width=3.9cm, height=1.75cm, scale only axis, at={(0,0)},
  xmin=2005, xmax=2025, ymin=0, ymax=496, axis x line*=bottom, axis y line=none,
  xtick={2005,2015,2025}, x tick label style={/pgf/number format/1000 sep={}, font=\small},
  title={Twitter}, title style={font=\large\bfseries, yshift=-1mm}]
\addplot[fill=gm@col@numbers!45!gmSurface, draw=gm@col@numbers, line width=1pt] coordinates {(2005,1) (2006,1) (2007,8) (2008,10) (2009,31) (2010,22) (2011,48) (2012,60) (2013,89) (2014,169) (2015,148) (2016,228) (2017,318) (2018,324) (2019,418) (2020,459) (2021,324) (2022,348) (2023,386) (2024,373) (2025,175)} \closedcycle;
\node[font=\small\bfseries, text=gmInk, anchor=south] at (axis cs:2020,459) {2020};
\end{axis}
\node[font=\small, text=gmInk2, anchor=north west] at (0,-.5) {3,940 memes};
\end{scope}
\begin{scope}[xshift=9.50cm, yshift=-3.25cm]
\begin{axis}[gm axis, width=3.9cm, height=1.75cm, scale only axis, at={(0,0)},
  xmin=2005, xmax=2025, ymin=0, ymax=496, axis x line*=bottom, axis y line=none,
  xtick={2005,2015,2025}, x tick label style={/pgf/number format/1000 sep={}, font=\small},
  title={TikTok}, title style={font=\large\bfseries, yshift=-1mm}]
\addplot[fill=gm@col@numbers!45!gmSurface, draw=gm@col@numbers, line width=1pt] coordinates {(2005,1) (2006,0) (2007,1) (2008,0) (2009,1) (2010,0) (2011,0) (2012,0) (2013,1) (2014,1) (2015,0) (2016,2) (2017,5) (2018,25) (2019,87) (2020,219) (2021,244) (2022,360) (2023,344) (2024,273) (2025,211)} \closedcycle;
\node[font=\small\bfseries, text=gmInk, anchor=south] at (axis cs:2022,360) {2022};
\end{axis}
\node[font=\small, text=gmInk2, anchor=north west] at (0,-.5) {1,775 memes};
\end{scope}
\end{tikzpicture}
}
  \vskip3mm
  {\small\color{gmInk2}Peaks: 4chan 2010, Tumblr 2014, YouTube and Reddit 2019, Twitter 2020, TikTok 2022.}
  \note{
  \begin{itemize}
    \item From the origin field, canonicalised (``Twitter'', ``twitter.com'', ``X'' are one
      platform), counted per year.
    \item The eras of Internet culture appear by themselves: 4chan peaks in 2010, Tumblr in 2014,
      YouTube and Reddit in 2019, Twitter in 2020 (by far the biggest: 3,940 memes since 2005),
      TikTok in 2022 --- from nothing before 2018.
    \item Each panel has the same scale, so the heights compare directly.
  \end{itemize}}
\end{frame}

\begin{frame}{What the graph talks about most}
  \begin{columns}[t]
    \column{.5\textwidth}
      {\scriptsize\color{gmInk2}Wikidata items, by entries that name them}\par\vskip1mm
      % generated by scripts/figures.py — do not edit
\begin{tikzpicture}
\begin{axis}[gm hbar, scale only axis, width=.8\linewidth, height=5.0cm, xmax=4768, symbolic y coords={{Twitter},{TikTok},{image macro},{YouTube},{catchphrase},{video game},{parody},{viral video},{song},{Reddit}},
  point meta=explicit symbolic, nodes near coords={\pgfplotspointmeta}]
\addplot[fill=kWikidata] coordinates {(3725,{Twitter}) [3,725] (3277,{TikTok}) [3,277] (2529,{image macro}) [2,529] (2260,{YouTube}) [2,260] (2138,{catchphrase}) [2,138] (1515,{video game}) [1,515] (1485,{parody}) [1,485] (1473,{viral video}) [1,473] (1429,{song}) [1,429] (1297,{Reddit}) [1,297]};
\end{axis}
\end{tikzpicture}

    \column{.46\textwidth}
      {\scriptsize\color{gmInk2}templates, by entries that use them}\par\vskip1mm
      % generated by scripts/figures.py — do not edit
\begin{tikzpicture}
\begin{axis}[gm hbar, scale only axis, width=.8\linewidth, height=4.3cm, xmax=14, symbolic y coords={{Trollface},{Drake Hotline Bling},{Woman Yelling At Cat},{Soyboy Vs Yes Chad},{Roll Safe Think About It},{Buff Doge vs. Cheems},{Doge},{chudjak}},
  point meta=explicit symbolic, nodes near coords={\pgfplotspointmeta}]
\addplot[fill=kTemplate] coordinates {(10,{Trollface}) [10] (9,{Drake Hotline Bling}) [9] (9,{Woman Yelling At Cat}) [9] (8,{Soyboy Vs Yes Chad}) [8] (7,{Roll Safe Think About It}) [7] (6,{Buff Doge vs. Cheems}) [6] (6,{Doge}) [6] (6,{chudjak}) [6]};
\end{axis}
\end{tikzpicture}

  \end{columns}
  \note{
  \begin{itemize}
    \item Most-named Wikidata items: the platforms (Twitter, TikTok, YouTube, Reddit, X,
      Instagram) and meme formats (image macro, catchphrase, parody, viral video, song, video game).
    \item Most-shared templates: Trollface, Drake Hotline Bling, Woman Yelling at a Cat, Soyboy vs.
      Yes Chad, Roll Safe --- and Buff Doge vs.\ Cheems and Doge: our dog again.
    \item A template shared by several entries is a bridge between memes that look alike.
  \end{itemize}}
\end{frame}

\begin{frame}{Doge's family}
  \centering
  \resizebox{\linewidth}{!}{% generated by scripts/figures.py — do not edit
\begin{tikzpicture}[x=2.75cm, y=1cm, fam/.style={rounded corners=1.2mm, inner xsep=1.4mm, inner ysep=.9mm, font=\fontsize{6.2}{7}\selectfont, text width=2.45cm, align=center}]
\node[fam, fill=kFrame!14!gmSurface, draw=kFrame!45!gmSurface] (n66d62ed870) at (0,-4.65) {Memes};
\node[fam, fill=kFrame!14!gmSurface, draw=kFrame!45!gmSurface] (n46f252304f) at (1,-4.65) {Image Macros};
\node[fam, fill=kFrame!14!gmSurface, draw=kFrame!45!gmSurface] (n7ec47ee030) at (2,-4.65) {Interior Monologue Captioning};
\node[fam, fill=kFrame, text=white, font=\scriptsize\bfseries] (n136de6de72) at (3,-4.65) {Doge};
\node[fam, fill=gmSurface, draw=kFrame!45!gmSurface] (nbb2493a99d) at (4,0.00) {Doge 2 / Caesar};
\node[fam, fill=gmSurface, draw=kFrame!45!gmSurface] (na76d99d68e) at (4,-4.07) {Ironic Doge Memes};
\node[fam, fill=gmSurface, draw=kFrame!45!gmSurface] (na111e6e580) at (5,-2.42) {Dogelore};
\node[fam, fill=gmSurface, draw=kFrame!45!gmSurface] (nee7be4b5bf) at (6,-1.32) {Cheems};
\node[fam, fill=gmSurface, draw=kFrame!45!gmSurface] (nef7a91ad8a) at (7,-0.66) {Bonk (Cheems)};
\node[fam, fill=gmSurface, draw=kFrame!45!gmSurface] (nfc77a9153f) at (7,-1.32) {Drog / Wo De Dāo Dùn / Cheems Frog};
\node[fam, fill=gmSurface, draw=kFrame!45!gmSurface] (n69d716db40) at (7,-1.98) {Slav Cat / Slav Cheebs};
\node[fam, fill=gmSurface, draw=kFrame!45!gmSurface] (nb6d019c33e) at (6,-2.64) {Quoge};
\node[fam, fill=gmSurface, draw=kFrame!45!gmSurface] (n49647443a6) at (6,-3.30) {Uncle Murphy};
\node[fam, fill=gmSurface, draw=kFrame!45!gmSurface] (n4bed7e7c52) at (5,-3.96) {Swole Doge};
\node[fam, fill=gmSurface, draw=kFrame!45!gmSurface] (n8aaa2bae9d) at (6,-3.96) {Swole Doge vs. Cheems};
\node[fam, fill=gmSurface, draw=kFrame!45!gmSurface] (nb8837ef101) at (7,-3.96) {Jojo Doge vs. Cheems};
\node[fam, fill=gmSurface, draw=kFrame!45!gmSurface] (n28301de2fa) at (5,-4.62) {There Is No Meme / I Lied to You, Take Off Your Clothes};
\node[fam, fill=gmSurface, draw=kFrame!45!gmSurface] (n659dd20625) at (6,-4.62) {I Lied, I Don't Have Netflix};
\node[fam, fill=gmSurface, draw=kFrame!45!gmSurface] (ne4f410dd25) at (5,-5.28) {Trapped Doge};
\node[fam, fill=gmSurface, draw=kFrame!45!gmSurface] (n70a85bef65) at (4,-6.60) {Shiba Inus / Shibes};
\node[fam, fill=gmSurface, draw=kFrame!45!gmSurface] (n37741e422f) at (5,-5.94) {Mein Scheiss Hund Gniesbert};
\node[fam, fill=gmSurface, draw=kFrame!45!gmSurface] (n552314226d) at (5,-6.60) {Shiba Inu Barking in an Office Meeting};
\node[fam, fill=gmSurface, draw=kFrame!45!gmSurface] (n9c41f2470c) at (5,-7.26) {Shiba Inu Dancing};
\node[fam, fill=gmSurface, draw=kFrame!45!gmSurface] (nbd8f531d89) at (4,-7.92) {The Death Of Kabosu / Doge};
\draw[kFrame!50!gmSurface, line width=.5pt, -{Stealth[length=1.3mm]}] (nef7a91ad8a.west) -- (nee7be4b5bf.east);
\draw[kFrame!50!gmSurface, line width=.5pt, -{Stealth[length=1.3mm]}] (nee7be4b5bf.west) -- (na111e6e580.east);
\draw[kFrame!50!gmSurface, line width=.5pt, -{Stealth[length=1.3mm]}] (na111e6e580.west) -- (na76d99d68e.east);
\draw[kFrame!50!gmSurface, line width=.5pt, -{Stealth[length=1.3mm]}] (na76d99d68e.west) -- (n136de6de72.east);
\draw[kFrame!50!gmSurface, line width=.5pt, -{Stealth[length=1.3mm]}] (nbb2493a99d.west) -- (n136de6de72.east);
\draw[kFrame!50!gmSurface, line width=.5pt, -{Stealth[length=1.3mm]}] (nfc77a9153f.west) -- (nee7be4b5bf.east);
\draw[kFrame!50!gmSurface, line width=.5pt, -{Stealth[length=1.3mm]}] (nbd8f531d89.west) -- (n136de6de72.east);
\draw[kFrame!50!gmSurface, line width=.5pt, -{Stealth[length=1.3mm]}] (n659dd20625.west) -- (n28301de2fa.east);
\draw[kFrame!50!gmSurface, line width=.5pt, -{Stealth[length=1.3mm]}] (n28301de2fa.west) -- (na76d99d68e.east);
\draw[kFrame!50!gmSurface, line width=.5pt, -{Stealth[length=1.3mm]}] (nb8837ef101.west) -- (n8aaa2bae9d.east);
\draw[kFrame!50!gmSurface, line width=.5pt, -{Stealth[length=1.3mm]}] (n8aaa2bae9d.west) -- (n4bed7e7c52.east);
\draw[kFrame!50!gmSurface, line width=.5pt, -{Stealth[length=1.3mm]}] (n4bed7e7c52.west) -- (na76d99d68e.east);
\draw[kFrame!50!gmSurface, line width=.5pt, -{Stealth[length=1.3mm]}] (n37741e422f.west) -- (n70a85bef65.east);
\draw[kFrame!50!gmSurface, line width=.5pt, -{Stealth[length=1.3mm]}] (n70a85bef65.west) -- (n136de6de72.east);
\draw[kFrame!50!gmSurface, line width=.5pt, -{Stealth[length=1.3mm]}] (nb6d019c33e.west) -- (na111e6e580.east);
\draw[kFrame!50!gmSurface, line width=.5pt, -{Stealth[length=1.3mm]}] (n552314226d.west) -- (n70a85bef65.east);
\draw[kFrame!50!gmSurface, line width=.5pt, -{Stealth[length=1.3mm]}] (n9c41f2470c.west) -- (n70a85bef65.east);
\draw[kFrame!50!gmSurface, line width=.5pt, -{Stealth[length=1.3mm]}] (n69d716db40.west) -- (nee7be4b5bf.east);
\draw[kFrame!50!gmSurface, line width=.5pt, -{Stealth[length=1.3mm]}] (ne4f410dd25.west) -- (na76d99d68e.east);
\draw[kFrame!50!gmSurface, line width=.5pt, -{Stealth[length=1.3mm]}] (n49647443a6.west) -- (na111e6e580.east);
\draw[kFrame!50!gmSurface, line width=.5pt, -{Stealth[length=1.3mm]}] (n136de6de72.west) -- (n7ec47ee030.east);
\draw[kFrame!50!gmSurface, line width=.5pt, -{Stealth[length=1.3mm]}] (n7ec47ee030.west) -- (n46f252304f.east);
\draw[kFrame!50!gmSurface, line width=.5pt, -{Stealth[length=1.3mm]}] (n46f252304f.west) -- (n66d62ed870.east);
\end{tikzpicture}
}\par\vskip2mm
  {\small\color{gmInk2}\NSiblingPairs\ pairs of entries share a series --- only \NSiblingLinked\ of them (\NSiblingLinkedPct\%) link to each other on their pages.}
  \note{
  \begin{itemize}
    \item The series: KYM says ``part of a series on \dots''. Doge belongs to Interior Monologue
      Captioning, which belongs to Image Macros, then Memes. Below Doge: Ironic Doge Memes, then
      Dogelore, Cheems, Bonk Cheems --- four generations of one dog.
    \item From the series we derive siblings (7.0.0, Riccardo's request): two entries with the same
      parent. 679,392 pairs --- and only 3,602 of them link to each other on their pages. Series and
      links are different relations; the graph keeps both.
    \item The longest chain in the graph goes through Doge: Bonk Cheems $\to$ \dots\ $\to$ Memes,
      seven steps.
  \end{itemize}}
\end{frame}

\begin{frame}{Is the graph sound?}
  \begin{columns}[t]
    \column{.48\textwidth}
      {\scriptsize\bfseries\color{gm@col@numbers}clean}\par\vskip2mm
      \begin{tabular}{@{}rl@{}}
        \textbf{\NIsolated}       & isolated nodes\\
        \textbf{\NDuplicateEdges} & duplicate edges\\
        \textbf{0}                & cycles in a series\\
        \textbf{\NComponents}     & connected components\\
        \textbf{=}                & RDF re-derived independently\\
      \end{tabular}
    \column{.48\textwidth}
      {\scriptsize\bfseries\color{gmMuted}worth knowing}\par\vskip2mm
      \begin{tabular}{@{}rl@{}}
        \textbf{\NUnresolvedPct\%} & of series parents are not in the corpus\\
        \textbf{\NNoTypePct\%}     & of entries have no entry type\\
        \textbf{\FDupInGraph}       & entries held twice (fixed, ch.~2)\\
        \textbf{6,800}             & templates showing ``a man''\\
      \end{tabular}
  \end{columns}
  \vskip7mm
  {\small\color{gmInk2}Every build measures this, after it is published.}
  \note{
  \begin{itemize}
    \item kg/metrics.py runs after every publish, on the IMKG-comparable core: no isolated nodes,
      no duplicate edges, no cycle in any series chain, two connected components. And the RDF
      re-derivation through the mapping equals the published RDF.
    \item Worth knowing: 26\% of series parents (1,061 of 4,074) are pages not in the corpus (they
      exist as stubs); 18.8\% of entries have no entry type on KYM.
    \item Found preparing this talk, and fixed (chapter 2, gap 14): 813 entries are in this graph
      twice, at their old address and at the one KYM moved them to --- Doge among them. The 7.1.0
      graph holds each once.
    \item Generic image regions: ``a man'' in 6,800 templates --- the reason only three generic
      regions per template reach the graph.
  \end{itemize}}
\end{frame}

\begin{frame}{Try it}
  \begin{columns}[t]
    \column{.48\textwidth}
      {\scriptsize\bfseries\color{gm@col@numbers}Neo4j Browser}\par\vskip1mm
      {\ttfamily\scriptsize http://\textless host\textgreater:8080/browser/}\par\vskip2mm
      {\fontsize{6}{7.6}\selectfont\ttfamily
      MATCH (cur:KGPointer \{name:'current'\})\\
      MATCH (f:Frame \{build\_id: cur.build\_id,\\
      \hspace*{17mm}label: 'Doge'\})\\
      RETURN f,\\
      \hspace*{2mm}COLLECT \{ MATCH p = (f)-[:fromTitle|fromTags|\\
      \hspace*{6mm}fromAbout]->() RETURN p \} AS items,\\
      \hspace*{2mm}COLLECT \{ MATCH p = (f)-[:partOfSeries*1..5]->()\\
      \hspace*{6mm}RETURN p \} AS series}\par\vskip2mm
      {\scriptsize\color{gmInk2}colours: import \texttt{neo4j\_browser.grass}}
    \column{.48\textwidth}
      {\scriptsize\bfseries\color{gm@col@numbers}SPARQL, as IMKG would ask}\par\vskip1mm
      {\ttfamily\scriptsize http://\textless host\textgreater:8080/sparql/}\par\vskip2mm
      {\fontsize{6}{7.6}\selectfont\ttfamily
      PREFIX rdfs: <http://www.w3.org/2000/01/rdf-schema\#>\\
      PREFIX m4s:  <https://meme4.science/>\\
      SELECT ?sibling ?title WHERE \{\\
      \hspace*{2mm}<https://knowyourmeme.com/memes/doge>\\
      \hspace*{6mm}rdfs:seeAlso ?sibling .\\
      \hspace*{2mm}?sibling m4s:title ?title .\\
      \}}
  \end{columns}
  \note{
  \begin{itemize}
    \item If there is time: a live demo. Neo4j Browser with the shared stylesheet (frames blue,
      Wikidata orange, templates aqua, events yellow) --- the Doge query draws the graph of chapter 7.
    \item SPARQL: the query IMKG itself would write for an entry's siblings, rdfs:seeAlso, runs
      unchanged on MemeAtlas and returns the whole series.
    \item Both endpoints are live for the lab, behind the one open port.
  \end{itemize}}
\end{frame}

\gmzoomout{numbers}
